\section{Benchmark Construction}
\label{sec:benchmark-construction}

\paragraph{Unit of construction.}
HandoffBench targets state transfer at a fixed control boundary rather than
routing quality. Each episode fixes an authenticated upstream trace, source and
receiver roles, the boundary cut, a public action contract, deterministic mock
tools and user replies, and a machine-executable terminal predicate. The
evaluator-side boundary state is a set of typed atomic claims with exact trace
provenance and epistemic status (known, unknown, contradicted, or not applicable).
We include a claim only when changing or removing it can alter a legal next action
or terminal success. Goals, user constraints, verified business facts, ordinary
unresolved slots, scoped consent, and organizational policy checks are mutually
exclusive roles. Executable action preconditions remain in the public contract,
and tool-execution metadata remains in the trace; neither is duplicated as a
semantic claim for primary state $F_1$.

\paragraph{Generation and validation.}
We first authored workflow blueprints---roles, causal trace, authority boundary,
plausible alternatives, legal sequence, and first user-impacting action---then
instantiated synthetic dates, monetary values, statuses, scopes, attributes, and
opaque identifiers with deterministic generators. Families are the split unit:
entity replacement or paraphrase does not create a new family. The original pool
contained 200 families evenly divided among travel, commerce, procurement,
enterprise IT, and scheduling.

Every label-free packet was independently annotated twice and locked before
comparison. On 990 claims, claim $F_1$ was .961 and exact action-sequence
agreement was 174/200; category and criticality agreement had Cohen's
$\kappa=.945$ and $.242$, respectively, while status, typed value, and provenance
agreed on 990/990 claims. Disagreement-only adjudication covered all 739 queued
entries and rejected 24 tasks. Separately, a reviewer blind to labels and
adjudication examined 39 agreement-only tasks flagged by static validity risks
and rejected all: 34 had ambiguous terminal actions, four had ambiguous sequence
order, and one lacked public terminal semantics.

We generated 63 new families to replace these disjoint exclusions, without
inspecting any candidate-model output, and double-annotated them anew. The
annotators agreed on all 63 action sequences and on identity, category, status,
value, and provenance for all 299 claims. Criticality agreement was 252/299
($\kappa=.693$); a third annotator resolved all 47 criticality disagreements,
rejecting no replacement. The final population is therefore 137 retained
originals plus 63 replacements, again 40 per domain (Table~\ref{tab:sealed-stats}).

\paragraph{Isolation and lineage.}
The two human roles received independently shuffled packets containing only the
trace cut, role boundary, public action schemas, and deterministic scripted
transitions. They did not receive source annotations, evaluator predicates,
development outputs, family blueprints, or one another's responses. Each
response was schema-validated and content-hashed before comparison. The
adjudicator saw only queued field or sequence disagreements; exact-agreement
items were not reopened. Likewise, the independent validity reviewer saw only
the 39 designated public packets, not A/B labels, queues, adjudication decisions,
or evaluator gold. This separation matters because annotator agreement alone
cannot show that a public contract identifies a unique terminal behavior.

Rejected originals remain in the audit lineage but are absent from the final
task files. All 63 replacements use new task and family identities and were
subjected to the same schema, provenance, grounding, mutation, and simulator
checks before inclusion. Sixteen replacements had exact agreement and required
no adjudication; the other 47 each contributed a criticality disagreement that
was resolved without changing public task evidence. A composite readiness
artifact binds original and replacement locks, agreement reports, adjudications,
the blind review, final task files, and final audit by SHA-256. This makes the
reported 200-task composition reproducible without exposing model outcomes,
because none existed during selection.

\begin{table}[t]
\centering
\small
\caption{Human validation, sealing, and static diagnostics for the final v3
split. Probe rates are not model-performance results.}
\label{tab:sealed-stats}
\begin{tabular}{@{}lr@{}}
\toprule
Statistic & Count / rate \\
\midrule
Sealed families (five domains) & 200 (40 each) \\
Retained originals / replacements & 137 / 63 \\
Original adjudication / blind-review rejects & 24 / 39 \\
Replacement disagreements / rejects & 47 / 0 \\
Unique public legal terminal sequence & 200 / 200 \\
Exact-copy oracle execution & 200 / 200 \\
Catalog-only / predicate-only success & 7.0\% / 6.0\% \\
Oracle sequence + enum-first success & 61 / 200 (30.5\%) \\
Exact development identity/trace collisions & 0 / 0 \\
Normalized topology flags vs. development & 69 \\
Lexical pairs at Jaccard $\geq .80$ & 0 \\
Candidate-model calls before sealing & 0 \\
\bottomrule
\end{tabular}
\end{table}

\paragraph{Final static audit and seal.}
All preregistered hard checks pass: every final task has one public legal terminal
sequence without access to the hidden predicate, grounded irreversible arguments,
valid provenance and impact labels, and successful deterministic oracle execution.
There are no development collisions in family ID, entity pool, seed, exact or
normalized trace, or exact action graph, and no trace bigram pair reaches .80.
The split is bound to seal \texttt{hb-v3.1-9671bf1ff2d5-20260721} and canonical
dataset hash \texttt{9671bf1ff2d5...d892a93}. No candidate-model call occurred
during construction, annotation, replacement, auditing, or sealing.

The leakage audit obtains 14/200 from a catalog-only heuristic and 12/200 from
predicate-only features. A stronger diagnostic, accurately called
\emph{oracle sequence plus enum-first}, is given the evaluator-side ordered
action-name sequence and guesses the first public enum value; it succeeds on
61/200 (30.5\%; Wilson 95\% CI [24.5, 37.2]). It therefore does not measure what
a receiver can infer from public action names alone. Exact copy succeeds on
200/200 and establishes simulator consistency only.

Sixty-nine final tasks share a normalized action-graph hash with development
tasks. The normalization erases action names, keys, and values, and can collapse
common ask/authorize/commit skeletons; this is a topology diagnostic rather than
evidence of semantic duplication. Conversely, its concentration means that we
do not claim a complete structural holdout. These diagnostics were reported
without post-hoc filtering. The split is frozen for the preregistered evaluation,
but confirmatory execution remains explicitly unauthorized and no confirmatory
result is reported here.
